\documentclass[10pt,a4paper]{article}
\usepackage[utf8]{inputenc}
\usepackage{amsmath, amsthm, amssymb}
\usepackage{geometry}
\usepackage{enumitem}
\usepackage{bm}

\geometry{
  a4paper,
  left=20mm,
  right=20mm,
  top=25mm,
  bottom=25mm,
}

\begin{document}
\section*{Problem}
\begin{enumerate}[label=\fbox{\arabic*}]
  \item % Problem 1
        Let $n$ be a positive integer. Let $\zeta = e^{\pi \sqrt{-1} / n}$ be a primitive $2n$-th root of unity. For integers $0 \le a, b \le 2n-1$, define a complex $2 \times 2$ matrix $X_{a, b}$ as follows.
        \[
          X_{a, b} = \begin{pmatrix} \zeta^a & 0 \\ 0 & \zeta^b \end{pmatrix}.
        \]
        Moreover, we define a complex $2 \times 2$ matrix $S$ as follows.
        \[
          S = \begin{pmatrix} 0 & 1 \\ 1 & 0 \end{pmatrix}.
        \]
        Let $G$ be the subgroup of $\operatorname{GL}_2(\mathbb{C})$ generated by $X_{a, b} \ (0 \le a, b \le 2n-1)$ and $S$.
        \begin{enumerate}[label=(\arabic*)]
          \item Calculate the number of elements of order $2$ in $G$.
          \item Let $H$ be the subgroup of $G$ generated by $\{g^2 \mid g \in G\}$. Calculate $[G : H]$.
          \item Calculate the number of subgroups $K$ of $G$ satisfying $[G : K] = 2$.
        \end{enumerate}

  \item % Problem 2
        Let $F$ be an element of the polynomial ring $\mathbb{C}[X, Y]$ over $\mathbb{C}$ with two variables $X, Y$. Let $R$ be the residue ring $\mathbb{C}[X, Y]/(F)$. Let $x, y$ be the residue classes of $X, Y$ in $R$, respectively, and $\mathfrak{m} = xR + yR$. Let $R_\mathfrak{m}$ be the localization of $R$ at $\mathfrak{m}$.
        \begin{enumerate}[label=(\arabic*)]
          \item Assume that $F$ is written as $F(X, Y) = X^a + Y^b$ for some integers $a, b$ greater than or equal to $2$. Prove that the ideal $\mathfrak{m} R_\mathfrak{m}$ is not principal.
          \item Assume that $F(0, 0) = 0$ and $\left( \left( \frac{\partial F}{\partial X} \right)(0, 0), \left( \frac{\partial F}{\partial Y} \right)(0, 0) \right) \neq (0, 0)$ are satisfied. Prove that $\mathfrak{m} R_\mathfrak{m}$ is a principal ideal.
        \end{enumerate}

  \item % Problem 3
        We set $S^2 = \{ (x_1, x_2, x_3) \in \mathbb{R}^3 \mid x_1^2 + x_2^2 + x_3^2 = 1 \}$. We define a subset $M$ of $S^2 \times S^2$ by
        \[
          M = \{ ((x_1, x_2, x_3), (y_1, y_2, y_3)) \in S^2 \times S^2 \mid x_1 y_1 + x_2 y_2 + x_3 y_3 = 0, \ x_1 + y_1 = 0 \}.
        \]
        \begin{enumerate}[label=(\arabic*)]
          \item Show that $M$ is a $C^\infty$-submanifold of $S^2 \times S^2$.
          \item Show that $M$ is orientable.
          \item We define a map $f : M \to S^2$ by
          \[
            f((x_1, x_2, x_3), (y_1, y_2, y_3)) = (y_1, y_2, y_3).
          \]
          Show that for any choice of an orientation on $M$, we have
          \[
            \int_M f^*\omega = 0
          \]
          for any $2$-form $\omega$ on $S^2$.
        \end{enumerate}

  \item % Problem 4
        Let $S^1 = \{ (x, y) \in \mathbb{R}^2 \mid x^2 + y^2 = 1 \}$. We fix a point $a \in S^1$. We consider an equivalence relation $\sim$ on $S^1 \times S^1$ generated by
        \[
          (a, p) \sim (p, a), \quad (p \in S^1).
        \]
        Let 
        \[
          X = (S^1 \times S^1) / \sim
        \]
        be the quotient space. Compute the homology groups of $X$ with integer coefficients.

  \item % Problem 5
        Let $f : [0, \infty) \to (-1, \infty)$ be a continuous function. Set
        \[
          g(x) = \sup_{y \in [x, x+1]} |f(y) - f(x)| \quad (x \ge 0)
        \]
        and suppose that
        \[
          \lim_{x \to \infty} f(x) = 0 \quad \text{and} \quad \int_0^\infty \{|f(x)| + g(x)\} dx < \infty
        \]
        are satisfied. Find the following limit:
        \[
          \lim_{n \to \infty} \prod_{k=1}^\infty \left\{ 1 + f\left(\frac{k}{n}\right) \right\}^{1/n}.
        \]

  \item % Problem 6
        Let $L^1([0, \infty))$ be the real Banach space of integrable functions over $[0, \infty)$, and for $f \in L^1([0, \infty))$, we define
        \[
          \|f\|_1 = \int_0^\infty |f(x)| dx.
        \]
        For $f \in L^1([0, \infty))$, we define $H_f : L^1([0, \infty)) \to L^1([0, \infty))$ by
        \[
          (H_f g)(x) = \int_0^\infty f(x+y)g(y)dy.
        \]
        \begin{enumerate}[label=(\arabic*)]
          \item Show that $\|H_f\| \le \|f\|_1$, where $\|H_f\|$ is the operator norm of $H_f$.
          \item Show that $H_f$ is a compact operator.
        \end{enumerate}

  \item % Problem 7
        Let $B := \{ x = (x_1, x_2) \in \mathbb{R}^2 \mid x_1^2 + x_2^2 < 1 \}$, $\bar{B} := \{ x = (x_1, x_2) \in \mathbb{R}^2 \mid x_1^2 + x_2^2 \le 1 \}$ and $\partial B$ be the boundary of $B$.
        \begin{enumerate}[label=(\arabic*)]
          \item Prove that if a real valued function $f \in C^1([0, 1])$ satisfies $f(1) = 0$, then
          \[
            \int_0^1 f(r)^2 r dr \le c_1 \int_0^1 f'(r)^2 r dr
          \]
          holds, where $c_1 > 0$ is a constant independent of $f$.
          \item Let $f \in C^1(B) \cap C(\bar{B})$ be a real valued function such that $f = 0$ on $\partial B$ and all of the first derivatives extend to $\bar{B}$ continuously. Prove that the inequality
          \[
            \int_B f(x)^2 dx \le c_2 \int_B \left\{ \left( \frac{\partial}{\partial x_1} f(x) \right)^2 + \left( \frac{\partial}{\partial x_2} f(x) \right)^2 \right\} dx
          \]
          holds with a constant $c_2 > 0$ independent of $f$.
          \item Let $u \in C^2((0, \infty) \times B) \cap C([0, \infty) \times \bar{B})$ be a real valued function satisfying the following conditions.
          \begin{itemize}
            \item $u(t, x) = 0, \ (t \in (0, \infty), x \in \partial B)$.
            \item For $t > 0$, $\frac{\partial}{\partial t} u(t, \cdot)$, $\frac{\partial}{\partial x_1} u(t, \cdot)$, $\frac{\partial}{\partial x_2} u(t, \cdot)$, $\frac{\partial^2}{\partial x_1^2} u(t, \cdot)$, $\frac{\partial^2}{\partial x_2^2} u(t, \cdot)$ extend to $\bar{B}$ continuously.
          \end{itemize}
          Prove that if $u$ solves
          \[
            \frac{\partial}{\partial t} u(t, x) - \frac{\partial^2}{\partial x_1^2} u(t, x) - \frac{\partial^2}{\partial x_2^2} u(t, x) = 0, \quad (t \in (0, \infty), x \in B),
          \]
          then for all $t > 0$, the inequality
          \[
            \int_B u(t, x)^2 dx \le c_3 e^{-c_4 t}
          \]
          holds, where $c_3 > 0$ is a constant independent of $t$, and $c_4 > 0$ is a constant independent of $t$ and $u$.
        \end{enumerate}

  \item % Problem 8
        Let $0 < a < 1$. Consider the following system of ordinary differential equations for a function $(r(t), \theta(t)) \in (0, \infty) \times (\mathbb{R}/2\pi\mathbb{Z})$ of $t \in \mathbb{R}$:
        \[
          \dot{r} = r \sin \theta, \quad \dot{\theta} = \frac{1}{r^2} + \cos \theta - a,
        \]
        where the dot ``$\cdot$'' means the differentiation with respect to $t$. Answer the following questions.
        \begin{enumerate}[label=(\arabic*)]
          \item Find equilibria of this system, and determine their linear stability.
          \item Prove that if $(r(t), \theta(t))$ is a solution of this system, then $(r(-t), -\theta(-t))$ is also a solution.
          \item Prove that there exists a solution $(r(t), \theta(t))$ of this system satisfying the following two conditions for some equilibrium $(\hat{r}, \hat{\theta})$:
          \begin{enumerate}
            \item[(a)] $\lim_{t \to \infty} (r(t), \theta(t)) = \lim_{t \to -\infty} (r(t), \theta(t)) = (\hat{r}, \hat{\theta})$.
            \item[(b)] For any $t \in \mathbb{R}$, $(r(t), \theta(t)) \neq (\hat{r}, \hat{\theta})$.
          \end{enumerate}
        \end{enumerate}

  \item % Problem 9
        Let $n$ be a nonnegative integer and define $A = \{i \mid 0 \le i \le n\}$. Suppose $a_0, a_1, \dots, a_n$ is a sequence of elements of $A$ such that $i \le a_i$ for every $i \in A$. Let us define a binary relation $\sim$ over $A$ by $\sim\ = \{(i, a_i) \mid 0 \le i \le n\}$. We denote by $\sim^*$ the smallest equivalence relation containing $\sim$.
        
        We consider the following program $G$. In the program, $U, I, J, X_0, \dots, X_n$ are program variables and $Y \leftarrow e$ is a substitution that assigns the value of $e$ to $Y$.
        \begin{verbatim}
        U <- 1;
        while U != 0 do
            U <- 0;
            for I = 0 to n do
                J <- X_I;
                if X_I != X_J then U <- U + 1 endif;
                X_I <- X_J;
            done
        done
        \end{verbatim}
        \begin{enumerate}[label=(\arabic*)]
          \item Suppose the precondition $\bigwedge_{i=0}^n X_i = a_i$ holds before the execution of the program $G$. Show that the condition $\Phi_1$ given below is a loop invariant for the while loop in $G$.
          \[
            \Phi_1 \equiv \forall i \in A. (i \le X_i \land (X_i = i \Leftrightarrow i = a_i))
          \]
          \item Assuming the same precondition as in (1), give a condition $\Phi$ that satisfies the following properties.
          \begin{enumerate}
            \item[(a)] $\Phi$ is a loop invariant for the while loop in $G$.
            \item[(b)] $\Phi \land U = 0$ implies $\forall i, j \in A. (X_i = X_j \Leftrightarrow i \sim^* j)$.
          \end{enumerate}
        \end{enumerate}

\end{enumerate}

\newpage
\section*{Solution}
\begin{enumerate}[label=\fbox{\arabic*}]
  \item % Solution 1
        \begin{enumerate}[label=(\arabic*)]
          \item The group $G$ consists of diagonal matrices $X_{a,b}$ and anti-diagonal matrices $X_{a,b}S$. The total number of elements in $G$ is $2(2n)^2 = 8n^2$. An element $g \in G$ has order 2 if and only if $g^2 = I$ and $g \neq I$.
          
          For a diagonal matrix $g = X_{a,b}$, $g^2 = X_{2a,2b} = I$ implies $2a \equiv 0 \pmod{2n}$ and $2b \equiv 0 \pmod{2n}$. Thus $a, b \in \{0, n\}$. There are $2 \times 2 = 4$ such matrices. Excluding the identity $I$, there are $3$ diagonal elements of order 2.
          
          For an anti-diagonal matrix $g = X_{a,b}S = \begin{pmatrix} 0 & \zeta^a \\ \zeta^b & 0 \end{pmatrix}$, we have $g^2 = \zeta^{a+b} I$. Thus, $g^2 = I$ if and only if $a+b \equiv 0 \pmod{2n}$. For each choice of $a \in \{0, \dots, 2n-1\}$, there is a unique $b$ satisfying this condition. Thus, there are $2n$ anti-diagonal elements of order 2.
          
          In total, the number of elements of order 2 in $G$ is $2n + 3$.
          
          \item The subgroup $H$ is generated by squares of all elements in $G$. The square of a diagonal element is $X_{a,b}^2 = X_{2a,2b}$. The square of an anti-diagonal element is $(X_{a,b}S)^2 = \zeta^{a+b} I = X_{a+b, a+b}$. Since both types of squares are diagonal matrices, $H$ is a subgroup of the abelian group of diagonal matrices in $G$. 
          
          The group $H$ is generated by $X_{2,0}, X_{0,2}$ and $X_{1,1}$. Any element in $H$ is of the form $X_{2a, 2b} X_{k,k} = X_{2a+k, 2b+k}$. This implies that the diagonal entries $\zeta^c$ and $\zeta^d$ of any matrix in $H$ must have exponents of the same parity ($c \equiv d \pmod 2$). 
          Conversely, any diagonal matrix with $c \equiv d \pmod 2$ can be generated this way. There are $n \times n$ combinations where both are even, and $n \times n$ where both are odd, giving a total of $|H| = 2n^2$ elements.
          The index is therefore $[G : H] = \frac{|G|}{|H|} = \frac{8n^2}{2n^2} = 4$.
          
          \item Subgroups $K \subset G$ with $[G : K] = 2$ correspond to kernels of surjective homomorphisms $\phi : G \to \{1, -1\}$. Such homomorphisms must contain the commutator subgroup $[G, G]$ in their kernel.
          The commutators are $[X_{a,b}, X_{c,d}] = I$ and $[X_{a,b}, S] = X_{a-b, b-a}$. Thus, $[G, G]$ is generated by $X_{1, -1}$, meaning $\phi(X_{1, -1}) = 1$, which implies $\phi(X_{1,0}) = \phi(X_{0,1})$.
          Let $x = \phi(X_{1,0}) \in \{1, -1\}$ and $s = \phi(S) \in \{1, -1\}$. Then $\phi(X_{a,b}) = x^{a+b}$ and $\phi(X_{a,b}S) = x^{a+b}s$.
          The condition $x^{2n} = 1$ is trivially satisfied for $x \in \{1, -1\}$. Any choice of $x, s \in \{1, -1\}$ defines a valid homomorphism.
          Since $\phi$ must be surjective to yield $[G : K] = 2$, we cannot have $x=1$ and $s=1$. The remaining 3 pairs $(x, s) \in \{(1, -1), (-1, 1), (-1, -1)\}$ give distinct surjective homomorphisms. Thus, there are exactly $3$ such subgroups.
        \end{enumerate}

  \item % Solution 2
        \begin{enumerate}[label=(\arabic*)]
          \item In the local ring $R_\mathfrak{m}$, Nakayama's Lemma implies that the minimum number of generators of the maximal ideal $\mathfrak{m} R_\mathfrak{m}$ is equal to the dimension of the $\mathbb{C}$-vector space $\mathfrak{m} R_\mathfrak{m} / (\mathfrak{m} R_\mathfrak{m})^2$, which is isomorphic to $\mathfrak{m} / \mathfrak{m}^2$ as an $R/\mathfrak{m} \cong \mathbb{C}$ vector space.
          By definition, $\mathfrak{m} / \mathfrak{m}^2 \cong (X, Y) / ((X, Y)^2 + (F))$.
          Since $a, b \ge 2$, the polynomial $F(X, Y) = X^a + Y^b$ is contained entirely in $(X, Y)^2$. Therefore, the ideal generated by $F$ is absorbed into $(X, Y)^2$, yielding:
          \[
            (X, Y)^2 + (F) = (X, Y)^2.
          \]
          This means $\mathfrak{m} / \mathfrak{m}^2 \cong (X, Y) / (X, Y)^2$. The standard basis for this vector space over $\mathbb{C}$ is $\{X, Y\}$, so its dimension is 2. Since the minimum number of generators is 2, the ideal $\mathfrak{m} R_\mathfrak{m}$ cannot be generated by a single element, and thus it is not a principal ideal.

          \item If $F(0,0) = 0$ and the gradient at the origin is non-zero, then $F$ has no constant term, but it has a non-zero linear part. This means $F \in (X, Y)$ but $F \notin (X, Y)^2$. 
          In the quotient $(X, Y) / (X, Y)^2$, the element $F$ introduces exactly one non-trivial linear relation between $X$ and $Y$ over $\mathbb{C}$.
          This single independent constraint reduces the dimension of the space by 1:
          \[
            \dim_{\mathbb{C}} \frac{(X, Y)}{(X, Y)^2 + (F)} = 2 - 1 = 1.
          \]
          Because the dimension of $\mathfrak{m} / \mathfrak{m}^2$ is 1, Nakayama's Lemma guarantees that the maximal ideal $\mathfrak{m} R_\mathfrak{m}$ can be generated by exactly 1 element. Therefore, $\mathfrak{m} R_\mathfrak{m}$ is a principal ideal.
        \end{enumerate}

  \item % Solution 3
        \begin{enumerate}[label=(\arabic*)]
          \item Let $F : \mathbb{R}^6 \to \mathbb{R}^4$ be defined by $F(x, y) = (\|x\|^2 - 1, \|y\|^2 - 1, x \cdot y, x_1 + y_1)$. Then $M = F^{-1}(0, 0, 0, 0)$. 
          The Jacobian matrix of $F$ is given by
          \[
            dF = \begin{pmatrix} 2x^\top & 0 \\ 0 & 2y^\top \\ y^\top & x^\top \\ e_1^\top & e_1^\top \end{pmatrix},
          \]
          where $e_1 = (1, 0, 0)^\top$. We check if $dF$ has full rank (rank 4) for all $(x,y) \in M$ by finding the kernel of its transpose. If $dF^\top \lambda = 0$, then:
          \begin{align*}
            2\lambda_1 x + \lambda_3 y + \lambda_4 e_1 &= 0, \\
            2\lambda_2 y + \lambda_3 x + \lambda_4 e_1 &= 0.
          \end{align*}
          Taking the dot product of the first equation with $x$ and the second with $y$, and using the conditions $\|x\|=\|y\|=1, x\cdot y = 0$, we find $2\lambda_1 + \lambda_4 x_1 = 0$ and $2\lambda_2 + \lambda_4 y_1 = 0$.
          Taking the cross dot products (first with $y$, second with $x$) gives $\lambda_3 + \lambda_4 y_1 = 0$ and $\lambda_3 + \lambda_4 x_1 = 0$.
          These last two equations imply $\lambda_4 (x_1 - y_1) = 0$. Since $x_1 + y_1 = 0$ on $M$, this becomes $2\lambda_4 x_1 = 0$. 
          If $x_1 \neq 0$, then $\lambda_4 = 0$, which immediately forces $\lambda_1=\lambda_2=\lambda_3=0$.
          If $x_1 = 0$, then $y_1 = 0$. The equations simplify to $2\lambda_1 x + \lambda_3 y + \lambda_4 e_1 = 0$. Since $x_1=y_1=0$, taking the dot product with $e_1$ yields $\lambda_4 = 0$, which again implies all $\lambda_i = 0$.
          Thus, the Jacobian has full rank everywhere on $M$. By the Regular Level Set Theorem, $M$ is a $C^\infty$-submanifold of dimension $6 - 4 = 2$.

          \item $M$ is defined as the preimage of a regular value of a smooth map into $\mathbb{R}^4$. The normal bundle of $M$ in $\mathbb{R}^6$ is globally spanned by the gradient vectors of the four components of $F$, meaning the normal bundle is trivial. Because the ambient space $\mathbb{R}^6$ is orientable and $M$ has a trivial normal bundle, $M$ itself is orientable.

          \item The map $f: M \to S^2$ is defined by $f(x, y) = y$. We calculate the topological degree of $f$. The preimage $f^{-1}(y)$ for a generic $y \in S^2$ consists of points $x \in S^2$ such that $x \cdot y = 0$ and $x_1 = -y_1$.
          This geometrically represents the intersection of the plane $x \cdot y = 0$ (a great circle on $S^2$) and the plane $x_1 = -y_1$. 
          These two planes intersect the sphere if and only if the distance from the origin to the intersection line is less than or equal to 1. 
          Specifically, we need $x_2^2 + x_3^2 = 1 - y_1^2$ and $x_2 y_2 + x_3 y_3 = y_1^2$. By the Cauchy-Schwarz inequality, a solution exists if and only if $y_1^4 \le (1 - y_1^2)(y_2^2 + y_3^2) = (1 - y_1^2)^2$, which simplifies to $|y_1| \le \frac{1}{\sqrt{2}}$.
          If we choose a regular value $y \in S^2$ such that $|y_1| > \frac{1}{\sqrt{2}}$ (e.g., $y = (1, 0, 0)$), the preimage $f^{-1}(y)$ is empty.
          Since the degree of a smooth map between connected orientable manifolds of the same dimension is constant for all regular values, we must have $\operatorname{deg}(f) = 0$.
          By the properties of the pullback integral, $\int_M f^*\omega = \operatorname{deg}(f) \int_{S^2} \omega = 0$.
        \end{enumerate}

  \item % Solution 4
        The space $S^1 \times S^1$ is the torus $T^2$. The equivalence relation $\sim$ identifies the subspace $S^1 \times \{a\}$ with $\{a\} \times S^1$ via the natural identity map $(p, a) \mapsto (a, p)$.
        We can analyze the quotient space $X$ through its CW complex structure. The torus has a standard CW decomposition with one 0-cell $e^0$, two 1-cells $e_1^1$ and $e_2^1$, and one 2-cell $e^2$. The attaching map for $e^2$ corresponds to the word $e_1^1 e_2^1 (e_1^1)^{-1} (e_2^1)^{-1}$.
        
        The identification $(p, a) \sim (a, p)$ glues the two 1-cells $e_1^1$ and $e_2^1$ together. In the quotient space $X$, they become a single 1-cell $e^1$.
        The boundary of the 2-cell $e^2$ in the quotient space follows the word $e^1 e^1 (e^1)^{-1} (e^1)^{-1} = 1$. This means that the attaching map for the 2-cell is null-homotopic.
        Thus, the CW complex for $X$ consists of one 0-cell, one 1-cell, and one 2-cell, with all boundary maps in the cellular chain complex being identically zero:
        \[
          0 \to \mathbb{Z} \xrightarrow{0} \mathbb{Z} \xrightarrow{0} \mathbb{Z} \to 0.
        \]
        The homology groups $H_k(X) = \operatorname{Ker}(\partial_k) / \operatorname{Im}(\partial_{k+1})$ are therefore given by the chain groups themselves:
        \begin{align*}
          H_0(X) &\cong \mathbb{Z}, \\
          H_1(X) &\cong \mathbb{Z}, \\
          H_2(X) &\cong \mathbb{Z}, \\
          H_k(X) &= 0 \quad (k \ge 3).
        \end{align*}
        Geometrically, since the 2-cell is attached trivially along a point, $X$ is homotopy equivalent to the wedge sum of a circle and a sphere, $S^1 \vee S^2$, which yields the same homology groups.

  \item % Solution 5
        Let $P_n = \prod_{k=1}^\infty \left\{ 1 + f\left(\frac{k}{n}\right) \right\}^{1/n}$. We analyze the logarithm of this product:
        \[
          \log P_n = \frac{1}{n} \sum_{k=1}^\infty \log\left( 1 + f\left(\frac{k}{n}\right) \right).
        \]
        We define $h(x) = \log(1 + f(x))$. Since $f(x) \to 0$ as $x \to \infty$, there exists $M > 0$ such that for $x > M$, $|f(x)| \le \frac{1}{2}$, giving $|h(x)| = |\log(1+f(x))| \le C|f(x)|$ for some constant $C$. Since $f$ is continuous and $\int_0^\infty |f(x)| dx \le \int_0^\infty \{|f(x)| + g(x)\} dx < \infty$, the function $h(x)$ is Lebesgue integrable on $[0, \infty)$.
        
        We view $\log P_n$ as a Riemann sum for $\int_0^\infty h(x) dx$. We can bound the discretization error $E_n$:
        \[
          E_n = \left| \int_0^\infty h(x) dx - \frac{1}{n} \sum_{k=1}^\infty h\left(\frac{k}{n}\right) \right| \le \sum_{k=1}^\infty \int_{(k-1)/n}^{k/n} \left| h(x) - h\left(\frac{k}{n}\right) \right| dx.
        \]
        Because $h'(u)$ is bounded near $0$, there is a constant $C'$ such that $|h(u) - h(v)| \le C'|u - v|$ globally for the range of $f$. 
        For $x \in [\frac{k-1}{n}, \frac{k}{n}]$, the distance $|\frac{k}{n} - x| \le \frac{1}{n}$. Assuming $n \ge 1$, we have $\frac{k}{n} \in [x, x+1]$. Therefore, $|f(x) - f(k/n)| \le \sup_{y \in [x, x+1]} |f(y) - f(x)| = g(x)$.
        To show this bound forces the error to $0$, we define a tighter local supremum $g_\delta(x) = \sup_{y \in [x, x+\delta]} |f(y) - f(x)|$.
        The error is bounded precisely by $C' \int_0^\infty g_{1/n}(x) dx$.
        Since $f$ is uniformly continuous on compact sets and approaches 0, $g_\delta(x) \to 0$ pointwise as $\delta \to 0$. Moreover, $g_\delta(x) \le g(x)$ for $\delta \le 1$. Because $g(x)$ is integrable, Lebesgue's Dominated Convergence Theorem ensures that $\int_0^\infty g_{1/n}(x) dx \to 0$ as $n \to \infty$.
        Thus, $E_n \to 0$, and the limit of the sum is exactly the integral:
        \[
          \lim_{n \to \infty} \log P_n = \int_0^\infty \log(1 + f(x)) dx.
        \]
        Exponentiating back, the limit is:
        \[
          \exp\left( \int_0^\infty \log(1 + f(x)) dx \right).
        \]

  \item % Solution 6
        \begin{enumerate}[label=(\arabic*)]
          \item Using the definition of the $L^1$ norm and Tonelli's theorem to swap the order of integration for non-negative functions, we get:
          \begin{align*}
            \|H_f g\|_1 &= \int_0^\infty \left| \int_0^\infty f(x+y) g(y) dy \right| dx \\
            &\le \int_0^\infty \int_0^\infty |f(x+y)| |g(y)| dy dx \\
            &= \int_0^\infty |g(y)| \left( \int_0^\infty |f(x+y)| dx \right) dy.
          \end{align*}
          The inner integral evaluates to $\int_y^\infty |f(u)| du$, which is strictly bounded above by the total integral over $[0, \infty)$, yielding $\|f\|_1$. Thus:
          \[
            \|H_f g\|_1 \le \int_0^\infty |g(y)| \|f\|_1 dy = \|f\|_1 \|g\|_1.
          \]
          This holds for all $g \in L^1$, proving that the operator norm satisfies $\|H_f\| \le \|f\|_1$.

          \item To show $H_f$ is compact, we first assume $f$ is a continuous function with compact support in $[0, R]$. Then $f(x+y) = 0$ whenever $x+y > R$. Consequently, for any $g$, $(H_f g)(x) = 0$ for $x > R$. The function $H_f g$ is supported in $[0, R]$.
          Let $B$ be the unit ball in $L^1([0, \infty))$. For $g \in B$, $|(H_f g)(x)| \le \sup |f| \|g\|_1 \le \sup |f|$, so the image $H_f(B)$ is uniformly bounded in $C([0, R])$.
          To show equicontinuity:
          \[
            |(H_f g)(x+h) - (H_f g)(x)| \le \int_0^R |f(x+h+y) - f(x+y)| |g(y)| dy \le \sup_u |f(u+h) - f(u)| \|g\|_1.
          \]
          Since continuous functions on compact supports are uniformly continuous, the supremum goes to zero as $h \to 0$ independently of $g$. Thus, $H_f(B)$ is equicontinuous. By the Arzelà-Ascoli theorem, $H_f(B)$ is relatively compact in $C([0, R])$. Since convergence in the supremum norm on a bounded interval implies convergence in the $L^1$ norm, $H_f(B)$ is relatively compact in $L^1([0, \infty))$, making $H_f$ a compact operator for such $f$.
          For an arbitrary $f \in L^1([0, \infty))$, since compactly supported continuous functions are dense in $L^1$, we can choose a sequence $f_n$ of such functions converging to $f$ in $L^1$. The inequality from part (1) gives $\|H_{f_n} - H_f\| = \|H_{f_n - f}\| \le \|f_n - f\|_1 \to 0$. Since the limit of compact operators under the operator norm topology is compact, $H_f$ is a compact operator for any $f \in L^1$.
        \end{enumerate}

  \item % Solution 7
        \begin{enumerate}[label=(\arabic*)]
          \item By the Fundamental Theorem of Calculus and the condition $f(1) = 0$, we write $f(r) = -\int_r^1 f'(s) ds$. Squaring both sides and applying the Cauchy-Schwarz inequality yields:
          \[
            f(r)^2 = \left( \int_r^1 f'(s) ds \right)^2 \le \left( \int_r^1 \frac{1}{s} ds \right) \left( \int_r^1 f'(s)^2 s ds \right) = (-\log r) \int_r^1 f'(s)^2 s ds.
          \]
          We can bound the rightmost integral by extending the lower bound to 0:
          \[
            f(r)^2 \le (-\log r) \int_0^1 f'(s)^2 s ds.
          \]
          Multiplying by $r$ and integrating both sides from $0$ to $1$:
          \[
            \int_0^1 f(r)^2 r dr \le \left( \int_0^1 (-r \log r) dr \right) \int_0^1 f'(s)^2 s ds.
          \]
          Evaluating the constant factor:
          \[
            \int_0^1 -r \log r dr = \left[ -\frac{r^2}{2} \log r + \frac{r^2}{4} \right]_0^1 = \frac{1}{4}.
          \]
          Thus, the inequality holds with $c_1 = \frac{1}{4}$.

          \item This is a specific instance of the Poincaré inequality on a unit disk. Expressing the Cartesian integral in polar coordinates $f(r, \theta)$:
          \[
            \int_B f(x)^2 dx = \int_0^{2\pi} \int_0^1 f(r, \theta)^2 r dr d\theta.
          \]
          For any fixed angle $\theta$, the radial function $r \mapsto f(r, \theta)$ satisfies the conditions of part (1). Applying the inequality to the inner integral:
          \[
            \int_0^1 f(r, \theta)^2 r dr \le c_1 \int_0^1 (\partial_r f(r, \theta))^2 r dr.
          \]
          In polar coordinates, the squared norm of the gradient is $|\nabla f|^2 = (\partial_r f)^2 + \frac{1}{r^2}(\partial_\theta f)^2$. Therefore, $(\partial_r f)^2 \le |\nabla f|^2$.
          Substituting this bound back and integrating over $\theta$:
          \[
            \int_B f(x)^2 dx \le c_1 \int_0^{2\pi} \int_0^1 (\partial_r f)^2 r dr d\theta \le c_1 \int_B |\nabla f|^2 dx.
          \]
          Thus, the inequality holds with $c_2 = c_1 = \frac{1}{4}$.

          \item The given equation is the standard heat equation $\frac{\partial u}{\partial t} = \Delta u$. We define the total thermal energy functional $E(t) = \int_B u(t, x)^2 dx$. Taking the time derivative:
          \[
            E'(t) = 2 \int_B u \frac{\partial u}{\partial t} dx = 2 \int_B u \Delta u dx.
          \]
          Using Green's first identity (integration by parts in higher dimensions), since the boundary condition specifies $u = 0$ on $\partial B$:
          \[
            E'(t) = -2 \int_B |\nabla u|^2 dx.
          \]
          From the Poincaré inequality derived in part (2), we know that $\int_B |\nabla u|^2 dx \ge \frac{1}{c_2} \int_B u^2 dx = \frac{1}{c_2} E(t)$. Substituting this gives a differential inequality:
          \[
            E'(t) \le -\frac{2}{c_2} E(t).
          \]
          Integrating this differential inequality using Gronwall's Lemma yields:
          \[
            E(t) \le E(0) e^{-\frac{2}{c_2} t}.
          \]
          The function $u(t, x)$ is given to be continuous up to $t=0$, meaning the initial energy $E(0) = \int_B u(0, x)^2 dx$ is finite and strictly depends on the initial state of $u$. Letting $c_3 = \int_B u(0, x)^2 dx$ (which is a constant independent of $t$) and $c_4 = \frac{2}{c_2} = 8$ (which is a universal constant independent of both $t$ and $u$), we obtain the final inequality:
          \[
            \int_B u(t, x)^2 dx \le c_3 e^{-c_4 t}.
          \]
        \end{enumerate}

  \item % Solution 8
        \begin{enumerate}[label=(\arabic*)]
          \item To find the equilibria, we set $\dot{r} = 0$ and $\dot{\theta} = 0$.
          From $\dot{r} = r \sin \theta = 0$, since $r \in (0, \infty)$, we must have $\sin \theta = 0 \implies \theta = 0$ or $\pi$.
          If $\theta = 0$, then $\cos \theta = 1$, and $\dot{\theta} = \frac{1}{r^2} + 1 - a = 0 \implies r^2 = \frac{1}{a-1}$. Because $0 < a < 1$, the right side is negative, which gives no real solution for $r$.
          If $\theta = \pi$, then $\cos \theta = -1$, and $\dot{\theta} = \frac{1}{r^2} - 1 - a = 0 \implies r^2 = \frac{1}{a+1}$. Since $a > 0$, we find a unique real positive solution $r = \frac{1}{\sqrt{a+1}}$. Thus, there is exactly one equilibrium point: $(\hat{r}, \hat{\theta}) = \left( \frac{1}{\sqrt{a+1}}, \pi \right)$.
          To determine linear stability, we evaluate the Jacobian matrix $J$:
          \[
            J = \begin{pmatrix} \frac{\partial \dot{r}}{\partial r} & \frac{\partial \dot{r}}{\partial \theta} \\ \frac{\partial \dot{\theta}}{\partial r} & \frac{\partial \dot{\theta}}{\partial \theta} \end{pmatrix} = \begin{pmatrix} \sin \theta & r \cos \theta \\ -\frac{2}{r^3} & -\sin \theta \end{pmatrix}.
          \]
          At the equilibrium $(\hat{r}, \pi)$:
          \[
            J = \begin{pmatrix} 0 & -\hat{r} \\ -\frac{2}{\hat{r}^3} & 0 \end{pmatrix}.
          \]
          The eigenvalues $\lambda$ are the roots of $\lambda^2 - \det(J) = 0$. Here $\det(J) = (-\hat{r})\left(-\frac{2}{\hat{r}^3}\right) = -\frac{2}{\hat{r}^2} < 0$. Therefore, $\lambda = \pm \frac{\sqrt{2}}{\hat{r}}$. Since the eigenvalues are purely real with opposite signs, the equilibrium is a saddle point, which is linearly unstable.

          \item Let $r_1(t) = r(-t)$ and $\theta_1(t) = -\theta(-t)$. Taking their derivatives:
          \begin{align*}
            \dot{r}_1(t) &= -\dot{r}(-t) = -r(-t) \sin(\theta(-t)) = -r_1(t) \sin(-\theta_1(t)) = r_1(t) \sin \theta_1(t), \\
            \dot{\theta}_1(t) &= \dot{\theta}(-t) = \frac{1}{r(-t)^2} + \cos(\theta(-t)) - a = \frac{1}{r_1(t)^2} + \cos(-\theta_1(t)) - a = \frac{1}{r_1(t)^2} + \cos \theta_1(t) - a.
          \end{align*}
          The transformed functions satisfy the exact same system of ODEs, confirming it is also a solution. This reveals the system is reversible under the time-reversal symmetry involution $S(r, \theta) = (r, -\theta)$.

          \item A solution satisfying both given conditions is by definition a non-trivial homoclinic orbit to the saddle equilibrium. 
          To prove its existence, we note that the exactness condition $M dr + N d\theta = 0$ is satisfied, implying the system admits a conserved quantity (a Hamiltonian-like function). Dividing the equations gives $\frac{dr}{d\theta} = \frac{r\sin\theta}{1/r^2 + \cos\theta - a}$. Rearranging terms allows us to integrate a potential function $F(r, \theta)$:
          \[
            F(r, \theta) = r(a - \cos \theta) + \frac{1}{r}.
          \]
          It's easy to verify that $\frac{d}{dt}F(r, \theta) = \frac{\partial F}{\partial r}\dot{r} + \frac{\partial F}{\partial \theta}\dot{\theta} = 0$.
          The value of this conserved quantity at the equilibrium is $F(\hat{r}, \pi) = 2\sqrt{a+1}$.
          A homoclinic orbit must reside entirely within this energy level set. Setting $F(r, \theta) = 2\sqrt{a+1}$ and solving the resulting quadratic equation $(a - \cos \theta)r^2 - 2\sqrt{a+1}r + 1 = 0$ for $r$ yields:
          \[
            r_\pm(\theta) = \frac{\sqrt{a+1} \pm \sqrt{1 + \cos \theta}}{a - \cos \theta}.
          \]
          We examine the branch $r_-(\theta) = \frac{\sqrt{a+1} - \sqrt{1 + \cos \theta}}{a - \cos \theta}$. By multiplying the numerator and denominator by the conjugate $\sqrt{a+1} + \sqrt{1+\cos \theta}$, the factor $(a - \cos \theta)$ cancels out beautifully:
          \[
            r_-(\theta) = \frac{(a+1) - (1+\cos \theta)}{(a-\cos \theta)(\sqrt{a+1} + \sqrt{1+\cos \theta})} = \frac{1}{\sqrt{a+1} + \sqrt{1+\cos \theta}}.
          \]
          Since $\cos \theta \ge -1$ and $a > 0$, this function $r_-(\theta)$ is strictly positive, uniformly bounded, and perfectly continuous for all $\theta \in [0, 2\pi]$. At $\theta = \pi$, it gracefully reaches $r_-(\pi) = \frac{1}{\sqrt{a+1}} = \hat{r}$, making it a smooth simple closed curve connected to the equilibrium point.
          Since the saddle point limits traversal time to infinity along its stable and unstable manifolds, the non-trivial trajectory that traces this continuous loop $r_-(\theta)$ starting from $\theta \to \pi$ requires an infinite amount of time $t \to \pm\infty$ to complete. This physically constitutes a strict homoclinic orbit, satisfying conditions (a) and (b).
        \end{enumerate}

  \item % Solution 9
        \begin{enumerate}[label=(\arabic*)]
          \item We prove that $\Phi_1$ is a loop invariant by showing it holds initially and is maintained after each execution of the inner \verb|for| loop.
          
          \textbf{Initialization:} Before the \verb|while| loop, the precondition asserts that $X_i = a_i$ for all $i \in A$. We are given that $i \le a_i$, which implies $i \le X_i$. Furthermore, $X_i = i$ is equivalent to $a_i = i$, which means $X_i = i \Leftrightarrow i = a_i$ is trivially satisfied. Thus, $\Phi_1$ holds initially.

          \textbf{Maintenance:} Assume $\Phi_1$ holds at the start of a \verb|while| loop iteration. We must show it still holds after the \verb|for| loop. We prove this step-by-step for each $I$-th iteration of the \verb|for| loop.
          Let $X_{\text{old}}$ and $X_{\text{new}}$ denote the states of the array $X$ immediately before and after the assignment $X_I \leftarrow X_J$ (where $J = X_{\text{old}}[I]$). Only the $I$-th element is modified, so $\Phi_1$ is trivially preserved for all $k \neq I$. For $k = I$, we must verify the two conditions in $\Phi_1$ for $X_{\text{new}}[I]$:
          \begin{itemize}
            \item Since $I \in A$, by the inductive hypothesis on $X_{\text{old}}$, we have $I \le X_{\text{old}}[I] = J$. Since $J \in A$, applying $\Phi_1$ again yields $J \le X_{\text{old}}[J]$. Thus, $I \le J \le X_{\text{old}}[J]$. Since $X_{\text{new}}[I] = X_{\text{old}}[J]$, we get $I \le X_{\text{new}}[I]$.
            \item We must show $X_{\text{new}}[I] = I \Leftrightarrow I = a_I$. 
            If $X_{\text{new}}[I] = I$, the inequality $I \le J \le X_{\text{new}}[I]$ squeezes $J$ to exactly $I$. Thus, $X_{\text{old}}[I] = I$. By $\Phi_1$ on $X_{\text{old}}$, this implies $I = a_I$.
            Conversely, if $I = a_I$, $\Phi_1$ on $X_{\text{old}}$ forces $X_{\text{old}}[I] = I$. Consequently, $J = I$, which means $X_{\text{new}}[I] = X_{\text{old}}[I] = I$.
          \end{itemize}
          This confirms that $\Phi_1$ is preserved across every step of the \verb|for| loop, making it a loop invariant for the \verb|while| loop.

          \item Let the condition $\Phi$ be defined as:
          \[
            \Phi \equiv \Phi_1 \land \forall i \in A. (i \sim^* X_i)
          \]
          \begin{enumerate}
            \item[(a)] \textbf{Proof that $\Phi$ is a loop invariant:}
            Since $\Phi_1$ is already proven to be a loop invariant, we only need to show that $i \sim^* X_i$ is maintained.
            Initially, $X_i = a_i$. By definition, $i \sim a_i$, which implies $i \sim^* a_i$. Thus, it holds before the loop.
            During the \verb|for| loop, the array is updated via $X_I \leftarrow X_{X_I}$. By the inductive hypothesis, $I \sim^* X_{\text{old}}[I]$ and $X_{\text{old}}[I] \sim^* X_{\text{old}}[X_{\text{old}}[I]] = X_{\text{new}}[I]$. Because $\sim^*$ is an equivalence relation, it is transitive, meaning $I \sim^* X_{\text{new}}[I]$. Therefore, $\Phi$ is a loop invariant.

            \item[(b)] \textbf{Proof of implication:}
            Assume $\Phi$ holds and $U = 0$. 
            The fact that $U = 0$ after the \verb|while| loop means that during the final pass of the \verb|for| loop, the condition $X_I \neq X_J$ was false for all $I \in A$. Since $J = X_I$, this means $X_I = X_{X_I}$ for all $I \in A$.
            Let $r_i = X_i$. Then the above condition states that $r_i = X_{r_i}$.
            From $\Phi_1$, we know $X_{r_i} = r_i \Leftrightarrow r_i = a_{r_i}$. Thus, $r_i = a_{r_i}$ for all $i \in A$, meaning $X_i$ is always a fixed point of the sequence defined by $a$.

            We claim that each equivalence class of $\sim^*$ contains exactly one such fixed point. 
            For any $x \in A$, we have $x \le a_x \le a_{a_x} \le \dots$. Since $A$ is a finite set, this monotonically increasing sequence must eventually stabilize at some fixed point $r = a_r$. Since $(x, a_x) \in \ \sim$, we have $x \sim^* r$. 
            If an equivalence class contained two distinct fixed points $r_1, r_2$, then $r_1 \sim^* r_2$. However, by definition, $\sim$ represents a directed forest where edges point from $i$ to $a_i$. A fixed point is the root of a tree in this forest. Two roots can only be in the same connected component ($\sim^*$) if one is an ancestor of the other (impossible for distinct roots) or if they share a common ancestor (impossible since roots have no outward edges besides self-loops). Thus, the fixed point is unique per equivalence class.

            From $\Phi$, we know $i \sim^* X_i$. Because $X_i$ is a fixed point, it must be the unique fixed point of the equivalence class containing $i$.
            Therefore, $X_i = X_j$ if and only if $i$ and $j$ share the same unique fixed point in their equivalence class, which is true if and only if $i \sim^* j$.
          \end{enumerate}
        \end{enumerate}

\end{enumerate}

\end{document}
